\section{External measurement validation and retained failures}
\label{sec:final-external-validation}

The following assays address whether a measurement detects its proposed
phenomenon. They do not directly estimate AI/Human classification accuracy.
Passing software or integrity checks and satisfying a scientific gate are
reported separately. Failed gates remain failed even when a subset of metrics
or a subsequent exploratory classifier appears promising.

An early, small Music Flamingo pilot on MIDI-grounded drum comparisons
failed the project's admission checks for dynamics and rhythm ranking.
The model was therefore not included in the quantitative detector, and
this limited trial does not support a broader assessment of its music
understanding capabilities. Detailed results and logs remain in the
standalone report\footnote{\path{artifacts/music_flamingo_midi_rhythm_probe_20260903/reports/FINAL_REPORT_EN.md}};
the thesis focuses on the subsequent measurement workflows.

\subsection{Dynamics, rhythm change, and section variety}

The controlled-change experiment tested variable gain, variable tempo, and
variable section lengths, retaining constant-gain and constant-tempo negative
controls. Its frozen family gate required a probability-of-superiority point
estimate of at least 0.75 and a bootstrap 95\% lower bound above 0.60.

\begin{table}[htbp]
\centering\small
\caption{External controlled-change family results. Superiority concerns
paired changes, not authorship classification.}
\label{tab:external-drp}
\begin{tabular}{lrrll}
\toprule
Family & Pairs & Superiority & 95\% interval & Gate\\
\midrule
Dynamics & 50 & 0.880 & [0.780, 0.960] & Pass\\
Rhythm change & 44 & 0.705 & [0.580, 0.818] & Fail\\
Section/bar variety & 44 & 0.523 & [0.466, 0.591] & Fail\\
\bottomrule
\end{tabular}
\end{table}

Dynamics therefore has evidence of sensitivity to amplitude variation, not
recovery of original MIDI velocities from mastered mixtures. Within R, tempo
entropy responded individually (superiority 0.917 on 36 observable pairs), but
that does not override the failed family gate. P was severely limited by
observability: some section endpoints were available on only one positive
pair. Reporting perfect ordering on that one pair would be misleading.

All three workflows were subsequently applied exploratorily to the existing
2,000-recording Human/Suno/HeartMuLa/ACE-Step collection. Pooled AUC was 0.875
for the spectral baseline, 0.892 with dynamics, and 0.914 with all three
additions. These historical predictive results do not retrospectively qualify
R or P as validated phenomenon measurements and are not Native30 results.

\subsection{SC: low-band spatial response and codec controls}

The low-band SC engineering gate used seven MedleyDB recordings held out
from the earlier 24-recording spatial-control development set. The admitted
candidate consists only of median absolute interchannel level difference
and median side-energy fraction in three bands below 6 kHz. The analysis
uses 44.1 kHz, a 4,096-sample periodic Hann window, hop 1,024, and no
centring or padding. Higher-band, coherence, phase-increment and IQR
diagnostics are not additional admitted predictors.

All 196 primary arithmetic checks and 28 eligible codec comparisons passed.
Across six descriptors and 28 codec/recording pairs, 166 of 168
nuisance-to-target ratios were at most 0.10; the largest was 0.138683.
Each descriptor met the frozen requirement that at least 90\% of its
ratios were at most 0.10 and none exceeded one. Independent replay
recomputed the checks and ratios, verified 435 raw and three gate products,
and checked 56 target-waveform hashes without importing the producer.

This threshold was selected after the development analysis, then frozen
before these seven recordings. Development had shown substantial AAC128
sensitivity in the 6--12 kHz side-energy descriptor: median codec-to-width
target ratio 0.621720, with six of 24 comparisons above one. The subsequent
low-band restriction is therefore prospective relative to the seven-record
gate, not a completely development-free choice. The seven tracks are
artist-proxy/track-disjoint from the preceding controls but reuse MedleyDB;
they are not a new independent AI/Human attribution population.

The result supports response to the specified spatial interventions relative
to the tested codecs. It does not establish universal codec invariance,
panning-motion understanding or authorship detection. The later Equal60
matched SC increments provide the separate predictive test and retain
negative transfer effects.
\begingroup\raggedright
Evidence under \path{artifacts/audio_phenomena_expansion_20260907/}:
\path{STEREO_LOWBAND_GATE_V2_RESULTS_EN.md},
\path{preregistration/stereo_lowband_gate_frozen_v2.json},
\path{results/stereo_lowband_gate_v2/gate_result.json},
\path{code/audit_stereo_lowband_gate_v2.py}, and
\path{audit/stereo_lowband_gate_parent_replay_v2.json}.\par
\endgroup

\subsection{Breath events: recall without usable precision}

The retained breath benchmark used 113 reviewed clips, twenty singer-prefix
groups, 141 merged high/medium reference events, and 38 clips with no reference
event. A known archive-identity conflict was excluded before this evaluation;
the conflicting candidates and exclusion evidence were retained.

Independent maximum-bipartite event matching required predicted duration at
most two seconds, onset error at most 200 ms, and positive temporal overlap.
It reproduced 137 true positives, 1,257 false positives, and four false
negatives. Precision was 0.098278, recall 0.971631, and F1 0.178502. Both
precision and F1 were below their required 0.70 threshold. All 38 no-event
clips had at least one detection, contributing 352 false-positive events.

Aggregate frame BA of 0.912017 was a diagnostic on a different unit and could
not substitute for event precision or F1. The integrity/recomputation audit
passed, but the measurement gate failed. B was not promoted to AI/Human
classification on this evidence.

\subsection{Vibrato: differential observability}

The VocalSet development assay opened 43 pairs (86 isolated monophonic
recordings from ten singers), while 47 evaluation pairs remained unopened.
Pairing used singer and repertoire/exercise/vowel identity, not time-aligned
or guaranteed equal-pitch performances. The subset had originally been selected
for breath annotation and was not a random sample of all VocalSet.

The frozen coarse pYIN workflow used 16-kHz audio, 1,024-sample frames,
160-sample hops, and a C2--C6 range. It produced a descriptor for 42/43 straight
recordings (97.67\%) but only 37/43 vibrato recordings (86.05\%). The 11.63
percentage-point availability gap exceeded the predeclared ten-point limit.
Among the 37 jointly observable pairs, differences were positive in 33, tied
in three, and negative in one. This favorable conditional ordering could not
bypass the observability failure. An independent reconciliation passed 883
checks; it did not establish continuous-F0 ground-truth fidelity. V remained
unadmitted, and the reserved pairs were not used to rescue the candidate.

\subsection{Envelope modulation: a near-threshold result remains a failure}

The reserved Q assay used 52 distinct-PCM notes from 26 instrument identities,
eight fixed interventions, and three views: original, MP3 128 kbit/s, and
AAC-LC requested at 64 kbit/s, yielding 1,248 measurements. These were reused
NSynth sample-library notes, not an independent performed-song domain.

The six descriptors measured fast modulation-power fraction and modulation
entropy in three carrier bands: 250--1,000, 1,000--3,000, and 3,000--7,000 Hz.
Carrier frequency and modulation frequency are different quantities. The
entropy intervention compared chirped amplitude modulation with 48-Hz
amplitude modulation at depth 0.6; the required note-level increase was 0.10,
with at least 90\% instrument-balanced success.

Of 146 prospective checks, 141 passed and five failed. High-carrier entropy
success was 88.462\% in all three views, below 90\%. Two additional AAC
nuisance checks failed for low- and high-carrier entropy. The maximum available
nuisance ratios remained below one; the failures concerned the required
fraction at or below 0.1, not catastrophic ratios above one. Passing the
remaining rate, depth, gain, polarity, and eligibility checks did not qualify
the full Q family. Selecting only successful descriptors after inspection
would not inherit the original gate's validity.

The consistency audit checked 6,659 products and recomputed all 146 decisions.
A separate DSP replay reproduced 1,875,744 spectrum bins with zero numerical
discrepancy using the same NumPy/SciPy libraries. This was not an
independent-library replication. Requested codec bitrates were not assumed
equal to observed stream rates, and the two recipes were not a bitrate-matched
codec ranking.

\subsection{Bicoherence: a passed controlled measurement gate}

BC measures consistent spectral coupling at the fixed 500/750/1,250-Hz
triplet. The reserved GuitarSet assay used 90 recordings and completed 1,440
measurements without failed or skipped cases: 630 float64, 270 float32
precision controls, and 540 codec-decoded measurements. All 72 scientific
checks were satisfied after independent numerical/accountability replay.

The closed-coupled minus independent-phase contrast was 0.863033 at $-6$ dB
and 0.936507 at 0 dB; all ninety recording contrasts were positive. Baseline
median absolute scalar errors were 0.000299 for MP3 128k and 0.000634 for
Opus 96k, with respective 95th percentiles 0.001563 and 0.005402. Eligibility
mask agreement was 1.0 under these baseline codec comparisons. These numbers
are controlled measurement responses and errors, not classification AUC.

An initial independent audit failed on a raw complex-sum mantissa comparison.
The corrected auditor made normalized column storage explicitly F-order and
asserted contiguity, preserving scientific tolerances, thresholds and producer
products. The original failure, diagnosis, correction authority and successful
replay were retained. Audit success and scientific gate satisfaction were not
treated as automatic classifier admission: Native30 transfer required a
separate fixed-view measurement protocol.

The controlled guitar backgrounds do not establish universal validity on
mastered mixtures. Arithmetic stereo downmix can cancel content, bicoherence
is nonlinear, and phase-stable linear tones can produce a high statistic.
The result supports a limited fixed-statistic measurement capability, not
proof of a generative mechanism or established AI/Human utility.

\subsection{References to retained experiments}

\begingroup\raggedright
The D/R/P protocol, paired measurements, and application results are in
\path{artifacts/dynamics_rhythm_change_detector_extension_20260903/}, including
\path{reports/REPORT_EN.md} and the retained controlled-change CSVs.
The remaining reports are under
\path{artifacts/audio_phenomena_expansion_20260907/}:
\path{BREATH_BENCHMARK_RESULTS_EN.md},
\path{VOCALSET_VIBRATO_DEVELOPMENT_RESULTS_EN.md},
\path{results/modulation_nsynth_gate_report_v1/RESULTS_EN.md}, and
\path{results/bc_reserved_verified_report_v2/REPORT_EN.md}.
These summaries point to their frozen implementations, protocols, measurements
and independent audits; they do not replace the underlying evidence.
\par\endgroup
